% =====================================================================
% VERSION ENSAE, REDUCTION SFE : chapitre 8, identification partielle et bornes de capital.
% Phase 3 : copie a texte constant de chapitres/10_identification_partielle.tex.
% Phase 4 (1er octobre 2026) : reecrit et condense ; formules, tableaux, figures, renvois,
% citation et nombres conserves. Le texte d'origine reste dans chapitres/ (version Institut)
% et dans l'historique Git.
% Lu par main_ensae.tex seul : ne jamais le faire lire par main_institut.tex,
% qui garde chapitres/.
% Scripts : 30 (bornes), 31 (comparaison des regimes), 33 (relatif/absolu), 34, 61, 67.
% =====================================================================
\chapter{Identification partielle et bornes de capital}
\label{chap:partielle}

La direction de la contagion n'étant pas identifiable (chapitre~\ref{ch:identifiabilite}), on peut
poser $W$ à dire d'expert et publier un besoin de capital ponctuel, ou caractériser l'ensemble des $W$
compatibles avec les contraintes posées et publier des extrema numériques. Nous suivons la seconde voie et la comparons à la
première. Les montants de ce chapitre valent pour le périmètre financier étudié, avec la
sévérité OpRisk, à titre illustratif.

\section{Modèle en couches}

Le socle, indépendant de $W$, regroupe la fréquence d'entrée surdispersée, l'accumulation par cause
commune et la sévérité GPD. Sans contagion ($W = 0$), le besoin de capital est de $5\,275$~M\euro{} dans le calage
de référence (sévérité OpRisk, fréquence de l'état conforme, gain de propagation $g=0{,}90$). Le préprint de cascade climatique \citep{ccrn2026} retire aussi la propagation dans une ablation, comme à $W=0$,
pour mesurer l'apport de la structure.

Toute matrice se décompose de façon unique en
\[
  W \;=\; S + A, \qquad S = \tfrac{1}{2}(W + W^{\top}), \qquad A = \tfrac{1}{2}(W - W^{\top}) .
\]
Ici $A_{jk}>0$ signifie que la contagion va plutôt de $j$ vers $k$, et $A_{jk}<0$ de $k$ vers $j$. Des
données de co-occurrence identifieraient $S$ au premier ordre sous amorce uniforme et faible intensité, mais pas $A$, la direction
(proposition~\ref{prop:nonid}). Faute d'un tel panel, nous posons $S$ à partir du barème d'expert
(chapitre~\ref{chap:cascade}) et ne faisons varier que $A$ : les bornes qui suivent sont donc
conditionnelles à cette partie symétrique posée. Nous définissons l'ensemble admissible par
\[
  \mathcal{W}(t) \;=\; \bigl\{\, S + A \;:\; A^{\top} = -A, \;\; |A_{jk}| \leq t\,S_{jk},
  \;\; \rho(S+A) < 1 \,\bigr\},
\]
où la contrainte $|A_{jk}| \leq S_{jk}$ découle des conditions de positivité $W_{jk}\geq 0$ et
$W_{kj}\geq 0$. Le paramètre $t\in[0,1]$ mesure ce que la donnée ignore de la direction, l'ignorance
totale $t=1$ correspondant aux sources publiques actuelles. Pour chacune des $10$ paires de piliers,
$A_{jk}/S_{jk}$ varie donc entre $-t$ et $t$ : les contraintes de signe définissent un pavé, intersecté avec la contrainte spectrale. Ses $2^{10} = 1\,024$ sommets fixent chaque paire à un extrême ; leur admissibilité est contrôlée avant évaluation.

% SOURCES-SCRIPTS: 31 30

\section{Capital sur l'ensemble admissible}

Nous appelons ici « bornes numériques » le minimum et le maximum calculés aux sommets admissibles, notés $V(1)$. Il ne s'agit pas de bornes globales démontrées sur $\mathcal W(1)$. À $t=1$, on obtient
\[
  \left[\min_{W\in V(1)}\widehat{\mathrm{SCR}}(W),\max_{W\in V(1)}\widehat{\mathrm{SCR}}(W)\right] \;=\; [\,6\,858 \;;\; 8\,697\,]~\text{M\euro{}} ,
\]
avec un facteur $1{,}27$ entre les deux bouts. La contagion ajoute entre $+30\,\%$ et
$+65\,\%$ au socle. Le point d'expert, $8\,110$~M\euro{}, n'est qu'un point de cet ensemble
(figure~\ref{fig:pave-sommets}).

\begin{figure}[H]
  \centering
  \figover{Z_pave_sommets.png}
  \caption{Le pavé admissible et le capital qu'il donne. (a)~Coupe sur deux des dix paires :
$A_{jk}/S_{jk}$ va de $-1$, tout de $k$ vers $j$, à $+1$, tout de $j$ vers $k$. (b)~Capital aux
$1\,024$ sommets et aux $3\,000$ points intérieurs, à $t = 1$ : aucun point intérieur ne sort des bornes.}
  \label{fig:pave-sommets}
\end{figure}

Les sommets ne bornent le capital que si les extrêmes y sont atteints, ce que rien ne garantit
\emph{a priori}. Augmenter $A_{jk}$ élève $W_{jk}$ mais abaisse $W_{kj}$, si bien que la
réponse n'est pas monotone en $A_{jk}$. Les $3\,000$ points intérieurs tirés uniformément restent tous
dans les bornes, en VaR comme en perte moyenne. Cette propriété reste vérifiée
numériquement, mais non démontrée. Ces extrema numériques sont des quantiles de $40\,000$ années
simulées, au bruit non mesuré. À indice de queue élevé, la VaR est peu sensible à la structure de
cascade, la queue annuelle étant dominée par une perte unique. La perte moyenne, bornée
dans $[849\,;1\,100]$~M\euro{} et moins bruitée, et l'ordre de remédiation en rendent mieux compte.

Ramener l'ignorance de $t=1$ à $t=0{,}5$, c'est-à-dire borner l'asymétrie à la moitié de la
co-occurrence, réduirait la largeur des bornes de $1\,839$ à $695$~M\euro{}, soit $-62\,\%$. Ce scénario donne l'ordre de grandeur de ce que le reporting de la section~\ref{ch:livrable} permettrait de gagner, sans le garantir.

% SOURCES-SCRIPTS: 30 61

\section{Comparaison d'une matrice posée et de l'identification partielle}

\begin{center}\footnotesize
\renewcommand{\arraystretch}{1.3}
\begin{tabular}{@{}p{4.1cm} p{4.5cm} p{5.6cm}@{}}
\toprule
 & \textbf{Régime A : on pose $W$} & \textbf{Régime B : modèle en couches} \\
\midrule
Sortie sur le capital & un point, $8\,110$~M\euro{} & socle $5\,275$ + bande $[6\,858\,;8\,697]$ \\
Part indépendante de la direction & aucune, tout repose sur le classeur & $79\,\%$ (socle $61\,\%$, plancher commun aux sommets $18\,\%$), à $S$ posé \\
Panel biaisé & déplace le besoin de capital de $777$~M\euro{} & la bande ne bouge pas si $S$ reste fixé \\
Priorité de remédiation & P1, sans marge d'erreur & tête P1/P4, ordre interne ouvert \\
Si aucun expert ne répond & plus de matrice complète & exploration conditionnelle au $S$ d'expert maintenu \\
\bottomrule
\end{tabular}
\end{center}

\vspace{2pt}
\noindent Faire varier la direction posée à dire d'expert, depuis un panel qui confirmerait le classeur
jusqu'à un panel qui dirait l'inverse au bord du pavé admissible, déplace le besoin de capital du
régime~A de $777$~M\euro{}, soit $10\,\%$. La priorité bascule alors entre trois piliers, sans signal
pour l'utilisateur. La
bande du régime~B ne bouge pas. Sur l'ensemble admissible, P5 n'arrive jamais en tête, P3 dans moins
de $10\,\%$ des cas, et P1 ($45\,\%$) et P4 ($36\,\%$) concentrent $81\,\%$ des matrices testées. Un panel d'experts servirait donc à départager
P1 et P4 et à resserrer la bande ; le questionnaire diffusé est resté sans réponse.

% SOURCES-SCRIPTS: 31

\section{Niveaux et rapports de capital}

\begin{center}\footnotesize
\renewcommand{\arraystretch}{1.25}
\begin{tabular}{@{}p{4.3cm} r r r r@{}}
\toprule
\textbf{Calibration} & \textbf{socle} & \textbf{bornes} &
\textbf{haute/socle} & \textbf{\% indépendant de la direction} \\
\midrule
Référence (OpRisk)   & $5\,275$  & $[6\,858\,;8\,697]$   & $1{,}65$ & $78{,}9\,\%$ \\
Sévérité PRC         & $239$     & $[336\,;424]$         & $1{,}77$ & $79{,}2\,\%$ \\
Fréquence $\times 3$ & $11\,709$ & $[14\,353\,;17\,405]$ & $1{,}49$ & $82{,}5\,\%$ \\
\bottomrule
\end{tabular}
\end{center}

\vspace{2pt}
\noindent La borne haute explorée varie d'un facteur $41$ entre les calibrations PRC et fréquence triplée ; il ne s'agit pas d'un rapport uniforme de tous les niveaux de capital. Les
rapports varient d'au plus $33\,\%$ de leur valeur moyenne, et l'ordre de remédiation place les mêmes
deux piliers, P1 et P4, en tête dans les trois calibrations.

Le rapport le moins stable est la largeur de bande rapportée au socle, qui varie de $32{,}9\,\%$ de sa valeur moyenne et
diminue lorsque la fréquence augmente, la queue annuelle étant alors moins dominée par une perte
unique. Ces rapports et le classement des piliers résistent au changement de source (figure~\ref{fig:partielle}).
% SOURCES-SCRIPTS: 33

\section{Échelle de repli sur la direction}
\label{sec:echelle-u}

On peut écrire la direction par l'excès de propagation propre à la source,
$\operatorname{logit}p_{ij}=\operatorname{logit}p_j+u_{ij}$ ($p_j$ mesure l'attractivité de la cible
$j$), puis limiter $u$ par paliers (crans), de $u$ libre au cran nul où $u=0$. Une direction de rang $1$ (une seule valeur singulière de $u$, figure~\ref{fig:echelle-repli}b) capte $74\,\%$ de la
variance de la direction posée et resserre les bornes de $25\,\%$, au prix d'une
hypothèse de structure. Au cran
nul, P1 et P4 deviennent indiscernables (gains de remédiation $1\,051$ contre $1\,052$~M\euro{}),
alors que $353$~M\euro{} les séparent avec la direction du classeur. C'est donc la direction qui
départage les deux piliers prioritaires.

\begin{figure}[H]\centering
\figover{Z5_echelle_repli.png}
\caption{Échelle de repli sur la direction. (a)~Les bornes du capital se resserrent quand la
direction descend d'un cran : $u$ libre, rang~$1$, classeur, puis cran nul ($u=0$). (b)~Valeurs
singulières de la direction posée : le rang~$1$ en capte l'essentiel. (c)~Gains de remédiation par
pilier avec la direction du classeur et au cran nul.}
% SOURCES-SCRIPTS: 34
\label{fig:echelle-repli}
\end{figure}

\begin{figure}[H]
  \centering
  \figover{Z_identification_partielle.png}
  \caption{Identification partielle de la contagion dirigée. (a)~Les bornes se resserrent quand
l'ignorance directionnelle $t$ diminue. (b)~Le prior d'expert dans des bornes indépendantes du
choix de direction, conditionnelles à la partie symétrique posée. (c)~Priorité de remédiation (regret maximal). (d)~Dominance des
piliers sur l'ensemble admissible : P5 n'arrive jamais en tête, et P1 et P4 concentrent $81\,\%$ des
matrices testées.}
  \label{fig:partielle}
\end{figure}
% SOURCES-SCRIPTS: 30 67

% SOURCES-SCRIPTS: 34

Les capitaux des chapitres~\ref{chap:multietats} et~\ref{chap:resultats} reposent sur la marche auto-évitante ; les extrema ci-dessus, sur la cascade indépendante. Les deux ne se comparent pas directement.
